\section{模型的评价与推广}

\subsection{模型的优点}

本文模型首先实现了任务瞬时资源占用与小时能源结算的严格分离。任务只要与某小时存在正重叠，其 GPU 和瞬时 IT 功率便被计入容量审计；其实际重叠时长则用于计算 GPU-hour 和电量。该处理同时适应分钟级任务和小时级电力数据，避免了整小时近似造成的资源高估，也避免了小时平均负荷掩盖瞬时超限。

其次，算力、设施和能源边界被置于统一可审计框架内。AI IT 负荷由任务类型功率映射和真实调度重建，非 AI IT 负荷保持不可迁移，PUE 作用于总 IT 负荷；新能源被拆分为直接消纳、储能充电、外送和弃电，电网为储能充电产生的购电量也被计入成本和碳排放。由此，各项指标均可由任务与能源轨迹重新计算，而不依赖附件中的基准结果列。

再次，模型对时域末端和储能初始状态进行了明确处理。第 $2400$ 至第 $2405$ 小时只承接此前到达的弹性任务，第 $2406$ 小时仅用于终端结算；储能状态从第 $0$ 小时运行前的绝对初始 SOC 开始递推，并要求终端 SOC 不低于初始状态。无储能比较状态即使不可行，也不会被错误地传播为含储能模型或联合模型不可行，从逻辑上保持了不同方案之间的独立性。

最后，指标分母门禁、完整残差审计和启发式声明提高了结果的可解释性。网络时延、新能源利用率和服务质量只有在分母有效时才被计算，未定义维度不会以任意小常数强行进入评分。实际结果被明确标记为非 Pareto 方案且无最优性证书，避免把有限规则构造包装成严格优化结论。

\subsection{模型的不足}

模型的主要不足来自数学 formulation 与实际 implementation 之间的差距。需求预测采用固定 Huber 超参数，没有完成全部季节性候选和超参数组合的验证选模；任务调度没有实现有限冲突回溯，问题二没有执行单任务局部改进，问题三没有执行储能动作候选搜索，问题四也没有执行跨基准和情景共享的交替邻域搜索。因此，所得方案只保证已构造候选的全时域可行性，无法评价未探索邻域中是否存在成本、碳排放或峰值表现更优的方案。

储能策略的工程刻画仍较简化。模型没有计入电池循环退化、温度影响、自放电和寿命成本，实际规则型方案累计放电量为 $0.0$ MWh，因而没有形成可用于检验削峰和价格套利能力的有效充放电循环。含储能成本相对无储能状态增加 $304223.92909026146$ 元，只能说明当前规则型策略未实现经济性改善，不能据此否定储能优化本身的价值。

通信模型只使用确定的单向网络时延，没有考虑带宽拥塞、迁移数据量、传输能耗和通信费用。服务质量也采用平均等待的有限代理，而不是完整的等待时间归一化指标。对于数据规模更大、网络资源紧张或任务迁移代价显著的实际系统，这些省略可能改变任务迁移方向和多目标权衡关系。

情景分析同样受到固定任务方案的限制。问题四各情景保持联合基准的任务迁移方案不变，碳约束场景未执行碳定向搜索，因此未达到碳目标只说明声明启发式内没有得到达标候选。新能源波动情景虽使用固定随机种子保证可复现性，但单一固定种子不能充分描述随机新能源误差分布，后续仍需通过固定随机种子的重复仿真评价结果稳定性。

\subsection{模型的推广}

在算力调度方面，候选域方法可推广至多数据中心、边缘节点与云中心协同的任务分配问题。只需在候选域中加入节点可用性、数据驻留、可信执行环境或任务亲和性约束，即可处理具有隐私和合规要求的跨区域计算。对于允许抢占或分布式训练的任务，可将单一开工变量扩展为逐时执行份额，并增加任务连续性、检查点和同步通信约束。

在能源调度方面，区域独立平衡可进一步扩展为含输电线路的网络模型。通过引入节点相角、线路容量或直流潮流约束，可研究数据中心迁移与电力跨区传输之间的替代关系。储能模型还可加入退化成本、循环寿命和备用容量，使短期电价套利、削峰和长期设备损耗在同一目标中得到权衡。

在不确定性处理方面，新能源出力、任务到达和执行时长可被构造为情景树或分布鲁棒集合。灵敏度实验表明，运行成本对新能源水平较为敏感，因此新能源预测误差应优先纳入滚动时域决策。实际运行时，可在每个小时更新任务队列、负荷预测和新能源预测，并保留已开工任务的连续占用状态，从而形成可在线实施的模型预测控制框架\cite{rawlings2017model}。

在求解方法方面，可先以本文确定性贪心策略快速生成可行初始方案，再采用混合整数规划、大邻域搜索或分解协调方法改进任务和储能决策。若计算预算有限，任务候选可按时延、截止裕度和最低资源下界提前筛除；若需要多目标决策，则可使用 $\varepsilon$-约束法生成经过验证的非支配候选集\cite{miettinen1999nonlinear}。无论采用何种扩展，任务轨迹、能源轨迹和指标均应在完整时域上重新计算，并继续保留结构化不可用状态与约束残差审计。
